% Lines 4968-5465 from original attention_transformers_notes_ghost.tex
% Appendix: Case Studies and Benchmarks

%==============================================================================
\section{Case Studies and Benchmarks}
%==============================================================================

\subsection{Comparative Analysis of Attention Variants}

Let's implement a comprehensive benchmark comparing different attention mechanisms:

\begin{lstlisting}[language=Python, caption=Attention Mechanism Benchmark]
@dataclass(frozen=True)
class BenchmarkResult:
    method: str
    memory_usage: float  # MB
    computation_time: float  # ms
    accuracy: float  # If applicable
    parameters: int

class AttentionBenchmark:
    """Benchmark different attention mechanisms"""

    def __init__(
        self,
        seq_lengths: list[int],
        d_model: int,
        n_heads: int,
        batch_size: int,
        device: str = "cuda"
    ) -> None:
        self.seq_lengths = seq_lengths
        self.d_model = d_model
        self.n_heads = n_heads
        self.batch_size = batch_size
        self.device = device

        # Initialize attention mechanisms
        self.methods = {
            "standard": nn.MultiheadAttention(
                d_model, n_heads, batch_first=True
            ),
            "linear": LinearMultiHeadAttention(
                LinearAttentionConfig(d_model, n_heads)
            ),
            "sparse_local": self._create_sparse_attention("local"),
            "sparse_strided": self._create_sparse_attention("strided"),
        }

        # Move to device
        for method in self.methods.values():
            method.to(device)

    def _create_sparse_attention(self, pattern: str) -> nn.Module:
        """Create sparse attention module"""
        # Placeholder for sparse attention
        return nn.MultiheadAttention(
            self.d_model, self.n_heads, batch_first=True
        )

    def benchmark(self) -> pd.DataFrame:
        """Run comprehensive benchmark"""
        results = []

        for seq_len in self.seq_lengths:
            # Create dummy input
            x = torch.randn(
                self.batch_size, seq_len, self.d_model
            ).to(self.device)

            for method_name, method in self.methods.items():
                result = self._benchmark_single(
                    method_name, method, x
                )
                result['seq_length'] = seq_len
                results.append(result)

        return pd.DataFrame(results)

    def _benchmark_single(
        self,
        method_name: str,
        method: nn.Module,
        x: Tensor
    ) -> dict:
        """Benchmark single method"""

        # Warmup
        for _ in range(10):
            with torch.no_grad():
                _ = method(x, x, x)

        # Memory measurement
        torch.cuda.empty_cache()
        torch.cuda.reset_peak_memory_stats()

        start_mem = torch.cuda.memory_allocated() / 1024 / 1024  # MB

        # Time measurement
        start_time = time.time()

        with torch.no_grad():
            for _ in range(100):
                output = method(x, x, x)
                if isinstance(output, tuple):
                    output = output[0]

        end_time = time.time()
        peak_mem = torch.cuda.max_memory_allocated() / 1024 / 1024  # MB

        # Calculate metrics
        computation_time = (end_time - start_time) / 100 * 1000  # ms
        memory_usage = peak_mem - start_mem

        # Count parameters
        params = sum(p.numel() for p in method.parameters())

        return {
            'method': method_name,
            'memory_usage': memory_usage,
            'computation_time': computation_time,
            'parameters': params
        }

    def plot_results(self, df: pd.DataFrame) -> None:
        """Visualize benchmark results"""

        fig, axes = plt.subplots(2, 2, figsize=(15, 10))

        # Time vs sequence length
        ax = axes[0, 0]
        for method in df['method'].unique():
            method_df = df[df['method'] == method]
            ax.plot(
                method_df['seq_length'],
                method_df['computation_time'],
                marker='o',
                label=method
            )
        ax.set_xlabel('Sequence Length')
        ax.set_ylabel('Computation Time (ms)')
        ax.set_title('Computation Time vs Sequence Length')
        ax.legend()
        ax.grid(True)

        # Memory vs sequence length
        ax = axes[0, 1]
        for method in df['method'].unique():
            method_df = df[df['method'] == method]
            ax.plot(
                method_df['seq_length'],
                method_df['memory_usage'],
                marker='o',
                label=method
            )
        ax.set_xlabel('Sequence Length')
        ax.set_ylabel('Memory Usage (MB)')
        ax.set_title('Memory Usage vs Sequence Length')
        ax.legend()
        ax.grid(True)

        # Time complexity (log scale)
        ax = axes[1, 0]
        for method in df['method'].unique():
            method_df = df[df['method'] == method]
            ax.loglog(
                method_df['seq_length'],
                method_df['computation_time'],
                marker='o',
                label=method
            )
        ax.set_xlabel('Sequence Length (log)')
        ax.set_ylabel('Computation Time (log ms)')
        ax.set_title('Time Complexity Analysis')
        ax.legend()
        ax.grid(True)

        # Parameter efficiency
        ax = axes[1, 1]
        methods = df['method'].unique()
        params = [df[df['method'] == m]['parameters'].iloc[0] for m in methods]
        avg_time = [df[df['method'] == m]['computation_time'].mean() for m in methods]

        ax.scatter(params, avg_time, s=100)
        for i, method in enumerate(methods):
            ax.annotate(method, (params[i], avg_time[i]))

        ax.set_xlabel('Number of Parameters')
        ax.set_ylabel('Average Computation Time (ms)')
        ax.set_title('Parameter Efficiency')
        ax.grid(True)

        plt.tight_layout()
        plt.show()
\end{lstlisting}

\subsection{Real-World Applications}

\subsubsection{Document Retrieval with Transformers}

Implement a transformer-based document retrieval system:

\begin{lstlisting}[language=Python, caption=Document Retrieval System]
@dataclass(frozen=True)
class DocumentRetrieverConfig:
    max_doc_length: int = 512
    max_query_length: int = 64
    d_model: int = 768
    n_heads: int = 12
    n_layers: int = 6
    retrieval_method: str = "dot_product"  # "dot_product", "cosine", "learned"

class DocumentEncoder(nn.Module):
    """Encode documents into dense representations"""

    def __init__(self, config: DocumentRetrieverConfig) -> None:
        super().__init__()
        self.config = config

        # Token embeddings
        self.token_embed = nn.Embedding(30522, config.d_model)  # BERT vocab size
        self.pos_embed = nn.Embedding(config.max_doc_length, config.d_model)

        # Transformer encoder
        encoder_layer = nn.TransformerEncoderLayer(
            d_model=config.d_model,
            nhead=config.n_heads,
            dim_feedforward=config.d_model * 4,
            batch_first=True
        )
        self.encoder = nn.TransformerEncoder(
            encoder_layer,
            num_layers=config.n_layers
        )

        # Pooling strategy
        self.pooler = nn.Linear(config.d_model, config.d_model)
        self.activation = nn.Tanh()

    def forward(
        self,
        input_ids: Tensor,
        attention_mask: Tensor
    ) -> Tensor:
        batch_size, seq_len = input_ids.shape

        # Embeddings
        token_emb = self.token_embed(input_ids)
        pos_ids = torch.arange(seq_len, device=input_ids.device)
        pos_emb = self.pos_embed(pos_ids).unsqueeze(0)

        embeddings = token_emb + pos_emb

        # Transformer encoding
        # Convert attention mask to transformer format
        mask = attention_mask.float()
        mask = mask.masked_fill(mask == 0, float('-inf'))
        mask = mask.masked_fill(mask == 1, float(0.0))

        encoded = self.encoder(embeddings, src_key_padding_mask=mask)

        # Pool to get document representation
        # Use [CLS] token representation
        cls_repr = encoded[:, 0]
        pooled = self.activation(self.pooler(cls_repr))

        return pooled

class CrossEncoderRanker(nn.Module):
    """Cross-encoder for re-ranking documents"""

    def __init__(self, config: DocumentRetrieverConfig) -> None:
        super().__init__()
        self.config = config

        # Shared token embeddings
        self.token_embed = nn.Embedding(30522, config.d_model)
        self.segment_embed = nn.Embedding(2, config.d_model)
        self.pos_embed = nn.Embedding(
            config.max_doc_length + config.max_query_length,
            config.d_model
        )

        # Transformer encoder
        encoder_layer = nn.TransformerEncoderLayer(
            d_model=config.d_model,
            nhead=config.n_heads,
            dim_feedforward=config.d_model * 4,
            batch_first=True
        )
        self.encoder = nn.TransformerEncoder(
            encoder_layer,
            num_layers=config.n_layers
        )

        # Classification head
        self.classifier = nn.Sequential(
            nn.Linear(config.d_model, config.d_model),
            nn.Tanh(),
            nn.Dropout(0.1),
            nn.Linear(config.d_model, 1)
        )

    def forward(
        self,
        query_ids: Tensor,
        doc_ids: Tensor,
        query_mask: Tensor,
        doc_mask: Tensor
    ) -> Tensor:
        batch_size = query_ids.shape[0]

        # Concatenate query and document
        input_ids = torch.cat([query_ids, doc_ids], dim=1)
        attention_mask = torch.cat([query_mask, doc_mask], dim=1)

        # Create segment embeddings
        query_len = query_ids.shape[1]
        doc_len = doc_ids.shape[1]
        segment_ids = torch.cat([
            torch.zeros(batch_size, query_len, dtype=torch.long),
            torch.ones(batch_size, doc_len, dtype=torch.long)
        ], dim=1).to(input_ids.device)

        # Embeddings
        token_emb = self.token_embed(input_ids)
        segment_emb = self.segment_embed(segment_ids)
        pos_ids = torch.arange(
            input_ids.shape[1],
            device=input_ids.device
        ).unsqueeze(0)
        pos_emb = self.pos_embed(pos_ids)

        embeddings = token_emb + segment_emb + pos_emb

        # Encode
        mask = attention_mask.float()
        mask = mask.masked_fill(mask == 0, float('-inf'))
        mask = mask.masked_fill(mask == 1, float(0.0))

        encoded = self.encoder(embeddings, src_key_padding_mask=mask)

        # Classify relevance
        cls_repr = encoded[:, 0]
        relevance_score = self.classifier(cls_repr)

        return relevance_score.squeeze(-1)

class HybridRetriever:
    """Combine dense and sparse retrieval"""

    def __init__(
        self,
        doc_encoder: DocumentEncoder,
        cross_encoder: CrossEncoderRanker,
        index_type: str = "faiss"  # "faiss", "annoy", "custom"
    ) -> None:
        self.doc_encoder = doc_encoder
        self.cross_encoder = cross_encoder
        self.index_type = index_type

        # Document index
        self.doc_embeddings: Tensor | None = None
        self.doc_ids: list[str] = []

        # For FAISS indexing
        if index_type == "faiss":
            import faiss
            self.index = None

    def index_documents(
        self,
        documents: list[dict[str, Any]],
        batch_size: int = 32
    ) -> None:
        """Index documents for retrieval"""

        self.doc_encoder.eval()
        all_embeddings = []

        with torch.no_grad():
            for i in range(0, len(documents), batch_size):
                batch = documents[i:i + batch_size]

                # Tokenize (simplified)
                input_ids = torch.tensor([
                    doc['token_ids'] for doc in batch
                ])
                attention_mask = torch.tensor([
                    doc['attention_mask'] for doc in batch
                ])

                # Encode
                embeddings = self.doc_encoder(input_ids, attention_mask)
                all_embeddings.append(embeddings.cpu())

                # Store document IDs
                self.doc_ids.extend([doc['id'] for doc in batch])

        # Combine embeddings
        self.doc_embeddings = torch.cat(all_embeddings, dim=0)

        # Build index
        if self.index_type == "faiss":
            import faiss

            # Normalize embeddings for cosine similarity
            embeddings_np = self.doc_embeddings.numpy()
            faiss.normalize_L2(embeddings_np)

            # Build index
            d = embeddings_np.shape[1]
            self.index = faiss.IndexFlatIP(d)  # Inner product
            self.index.add(embeddings_np)

    def retrieve(
        self,
        queries: list[dict[str, Any]],
        k: int = 100,
        rerank_top_k: int = 10
    ) -> list[list[dict[str, Any]]]:
        """Retrieve and rerank documents"""

        self.doc_encoder.eval()
        self.cross_encoder.eval()

        results = []

        with torch.no_grad():
            for query in queries:
                # Encode query
                query_ids = torch.tensor([query['token_ids']])
                query_mask = torch.tensor([query['attention_mask']])

                query_emb = self.doc_encoder(query_ids, query_mask)

                # Initial retrieval
                if self.index_type == "faiss":
                    query_np = query_emb.cpu().numpy()
                    faiss.normalize_L2(query_np)

                    distances, indices = self.index.search(query_np, k)

                    # Get top-k documents
                    top_docs = []
                    for idx, score in zip(indices[0], distances[0]):
                        top_docs.append({
                            'id': self.doc_ids[idx],
                            'initial_score': float(score),
                            'index': idx
                        })

                # Re-rank top documents with cross-encoder
                if rerank_top_k > 0 and len(top_docs) > 0:
                    rerank_docs = top_docs[:rerank_top_k]

                    # Prepare batch for cross-encoder
                    doc_indices = [doc['index'] for doc in rerank_docs]

                    # Get document tokens (simplified)
                    doc_ids_batch = torch.stack([
                        torch.tensor(documents[idx]['token_ids'])
                        for idx in doc_indices
                    ])
                    doc_mask_batch = torch.stack([
                        torch.tensor(documents[idx]['attention_mask'])
                        for idx in doc_indices
                    ])

                    # Expand query for batch
                    query_ids_batch = query_ids.expand(len(rerank_docs), -1)
                    query_mask_batch = query_mask.expand(len(rerank_docs), -1)

                    # Get relevance scores
                    relevance_scores = self.cross_encoder(
                        query_ids_batch,
                        doc_ids_batch,
                        query_mask_batch,
                        doc_mask_batch
                    )

                    # Update scores
                    for i, doc in enumerate(rerank_docs):
                        doc['rerank_score'] = float(relevance_scores[i])

                    # Sort by rerank score
                    rerank_docs.sort(
                        key=lambda x: x['rerank_score'],
                        reverse=True
                    )

                    # Combine reranked and remaining
                    final_docs = rerank_docs + top_docs[rerank_top_k:]
                else:
                    final_docs = top_docs

                results.append(final_docs)

        return results
\end{lstlisting}
